\section{The actual canonical coefficient in geometric specialization}

\begin{lemma}[The original canonical-twisted arithmetic model]
\label{lem:kn-arithmetic-canonical-specialization-model}
\lean{EnvelopeProperty.Vanishing.ArithmeticCanonicalSpecializationGeometry.canonicalModel,EnvelopeProperty.Vanishing.ArithmeticCanonicalSpecializationGeometry.canonicalModel_isInvertible}
\leanok
\uses{thm:kn-arithmetic-global-canonical-fibre}
Let $p:\mathcal Y\to\operatorname{Spec}R$ be the original arithmetic family,
smooth of relative dimension $n$, and let $L$ be an original invertible
module on $\mathcal Y$.  The actual module
$M=\omega_{\mathcal Y/R}\otimes L$ is constructed on this same whole family
and is invertible.
\end{lemma}
\begin{proof}
Use the constructed whole relative top exterior sheaf and its actual
standard-smooth canonical atlas.  Tensor this original invertible module
with the same original $L$ using the native sheafified tensor product.
Native stability of invertibility under this tensor proves the claim.
\end{proof}

\begin{theorem}[The canonical coefficient on a prescribed Cartesian field square]
\label{thm:kn-arithmetic-canonical-prescribed-field-square}
\lean{EnvelopeProperty.Vanishing.ArithmeticCanonicalSpecializationGeometry.squareIso,EnvelopeProperty.Vanishing.ArithmeticCanonicalSpecializationGeometry.squareIso_fst,EnvelopeProperty.Vanishing.ArithmeticCanonicalSpecializationGeometry.squareIso_snd,EnvelopeProperty.Vanishing.ArithmeticCanonicalSpecializationGeometry.nativeFieldCanonicalTensorIso,EnvelopeProperty.Vanishing.ArithmeticCanonicalSpecializationGeometry.squareCanonicalTensorIso}
\leanok
\uses{lem:kn-arithmetic-canonical-specialization-model}
For an actual field map $\varphi:R\to k$ and an actual Cartesian square
with original maps $g:X\to\mathcal Y$ and $q:X\to\operatorname{Spec}k$ over
$p$ and $\operatorname{Spec}\varphi$, there is a constructed whole-sheaf
isomorphism
\[
  g^*M\ \simeq\ \omega_{X/k}\otimes g^*L.
\]
The canonical object is the existing native object for the original $q$,
and the second tensor factor is the pullback by the original $g$.
\end{theorem}
\begin{proof}
The actual Cartesian square supplies its own isomorphism from the original
$X$ to the native categorical pullback.  Its two original projection
equations identify the first projection with $g$ and the second structural
projection with $q$.  The coefficient map $\varphi$ supplies its literal
$R$-algebra structure on $k$, so the proved whole canonical model comparison
applies to this native field fibre.  Pull it back through the constructed
scheme isomorphism.  Native pullback composition identifies the original
model pullback, and the field-relative canonical transport retains the
actual second-projection triangle.  The final tensor comparison retains
the same original $g$ and $L$.  Every scheme and module comparison is
derived from the actual Cartesian square.
\end{proof}

\begin{lemma}[Actual canonical positive cohomology on the prescribed square]
\label{lem:kn-arithmetic-canonical-square-cohomology}
\lean{EnvelopeProperty.Vanishing.ArithmeticCanonicalSpecializationGeometry.nativeField_isProjectiveFactorization,EnvelopeProperty.Vanishing.ArithmeticCanonicalSpecializationCohomology.squareCanonicalTensor_subsingleton_H_iff}
\leanok
\uses{thm:kn-arithmetic-canonical-prescribed-field-square}
If the original $p$ has a genuine finite projective factorization, positive
cohomology of $\omega_{X/k}\otimes g^*L$ vanishes exactly when positive
cohomology of the same original $M$ pulled to the native categorical field
fibre vanishes.
\end{lemma}
\begin{proof}
The existing projective base-change construction supplies a projective
factorization for the actual native field fibre.  Invertibility of the
internally constructed $M$ supplies quasicoherence of its actual pullback.
The native positive cohomology comparison for a genuine scheme isomorphism
therefore applies to the actual Cartesian-square isomorphism.  Compare its
pulled coefficient with the prescribed native canonical tensor using the
preceding whole module isomorphism and the actual categorical sheaf
cohomology functor.  No cohomology comparison certificate is an input.
\end{proof}

\begin{theorem}[Canonical closed-fibre vanishing reaches the original ambient field]
\label{thm:kn-arithmetic-canonical-specialization}
\lean{EnvelopeProperty.Vanishing.ArithmeticCanonicalSpecializationCohomology.nativeCanonicalTensor_subsingleton_of_closed_of_injective,EnvelopeProperty.Vanishing.ArithmeticCanonicalSpecializationCohomology.squareCanonicalTensor_subsingleton_of_closed_of_injective}
\leanok
\uses{lem:kn-arithmetic-canonical-square-cohomology}
Assume $R$ is Noetherian and the original $p$ is smooth of relative dimension
$n$ with a genuine finite projective factorization.  Fix an actual maximal
ideal $P\subset R$, an original invertible $L$, and $j\geq0$.  If the actual
native closed fibre satisfies
\[
 H^{j+1}(\mathcal Y_{\kappa(P)},
   \omega_{\mathcal Y_{\kappa(P)}/\kappa(P)}\otimes L_{\kappa(P)})=0,
\]
then every prescribed injective map $\varphi:R\hookrightarrow F$ into an
ambient field, with its actual Cartesian square $(g,q)$ as above, satisfies
\[
  H^{j+1}(X,\omega_{X/F}\otimes g^*L)=0.
\]
\end{theorem}
\begin{proof}
Construct the same whole model $M$ internally and derive its invertibility
from the first lemma.  Original relative smoothness supplies flatness of
$p$.  The whole canonical field-fibre comparison converts the displayed
native closed canonical-twisted group to the actual closed pullback of $M$.
Apply the proved geometric cohomology specialization theorem for the same
original Noetherian flat projective family and invertible model.  Its
bounded flat ordered section complex and finite native homology produce an
original principal neighbourhood; injectivity of the prescribed field map
keeps its original defining coefficient nonzero in $F$.  Thus the actual
native ambient pullback of the same $M$ has zero positive cohomology.
The preceding Cartesian-square cohomology comparison gives the conclusion
on the original $X$, with its original $g$ and $q$.
\end{proof}

The closed-fibre premise here is the intermediate output of the actual
Deligne--Illusie and Witt-lift vanishing producer.  The characteristic-zero
entry composes that producer with this specialization theorem after
constructing the actual arithmetic model, thereby eliminating the
closed-fibre premise from the final Kodaira statement.

\begin{lemma}[The actual quotient and localization residue coefficients]
\label{lem:kn-arithmetic-canonical-quotient-residue}
\lean{EnvelopeProperty.Vanishing.ArithmeticCanonicalSpecializationQuotientGeometry.quotientResidueEquiv,EnvelopeProperty.Vanishing.ArithmeticCanonicalSpecializationQuotientGeometry.quotientResidueEquiv_mk,EnvelopeProperty.Vanishing.ArithmeticCanonicalSpecializationQuotientGeometry.residueFibreIso,EnvelopeProperty.Vanishing.ArithmeticCanonicalSpecializationQuotientGeometry.residueFibreIso_fst}
\leanok
For the actual maximal ideal $P$, the quotient field $R/P$ and the native
localization residue field $\kappa(P)$ have a constructed coefficient
isomorphism.  It takes every original residue class of $r\in R$ to the
actual image of $r$ in $\kappa(P)$.  Their actual categorical fibres have a
derived scheme isomorphism retaining the same original-family projection.
\end{lemma}
\begin{proof}
The actual algebra map $R/P\to\kappa(P)$ is bijective for a maximal ideal.
Its native ring equivalence gives a spectrum isomorphism in the opposite
direction.  The coefficient formula on all original residue classes proves
its actual structural triangle over $R$.  Native pullback lifts construct
the fibre map and its inverse; the actual two projection equations prove
their inverse identities.  The original-family projection is the same
first projection on both sides.
\end{proof}

\begin{theorem}[The literal quotient closed fibre supplies canonical specialization]
\label{thm:kn-arithmetic-canonical-quotient-specialization}
\lean{EnvelopeProperty.Vanishing.ArithmeticCanonicalSpecializationQuotient.nativeResidueCanonicalTensor_subsingleton_of_quotient,EnvelopeProperty.Vanishing.ArithmeticCanonicalSpecializationQuotient.squareCanonicalTensor_subsingleton_of_quotient_of_injective}
\leanok
\uses{lem:kn-arithmetic-canonical-quotient-residue,thm:kn-arithmetic-canonical-specialization}
In the preceding specialization theorem, the intermediate closed-fibre
vanishing may be supplied on the literal quotient fibre
$\mathcal Y\times_R\operatorname{Spec}(R/P)$ along the actual quotient map.
This is the actual finite field used by the produced Witt stage.  The same
conclusion holds on the prescribed original injective ambient field square.
\end{theorem}
\begin{proof}
The proved canonical comparison on the quotient field first converts its
native canonical-twisted group into the actual pullback of the same original
global model $M$.  Pullback along the derived fibre scheme isomorphism
preserves positive cohomology, using the actual projective factorization
over the quotient field and the internally derived quasicoherence of $M$.
Its original-family projection equation identifies the resulting coefficient
with the actual localization-residue pullback of that same $M$.
The proved canonical comparison within $\kappa(P)$ then supplies the native
closed-fibre hypothesis of the preceding theorem.  Apply that specialization
theorem to obtain the conclusion on the original ambient square.  Thus the
quotient and localization residue field are connected by actual ring,
scheme and model-pullback comparisons throughout.
\end{proof}
